\begin{example}
\label{example-noetherian-topology-on-spec}
\begin{reference}
Comment by Lukas Heger of November 12, 2020.
\end{reference}
Let $R$ be a ring and $X = \Spec(R)$. Since closed subsets of $X$ correspond to
radical ideals of $R$ (Lemma \ref{lemma-Zariski-topology}) we see that $X$ is a
Noetherian topological space if and only if we have ACC for radical ideals.
This holds if and only if every radical ideal is the radical of a finitely
generated ideal. For completeness, under ACC the set of radicals
of finitely generated subideals of a radical ideal $K$ has a maximal
member $\sqrt{F}$. Otherwise successive strict enlargements would
give an infinite ascending chain. For every $x\in K$ the ideal
$\sqrt{F+(x)}$ is another such member containing $\sqrt F$,
so maximality gives $x\in\sqrt F$ and hence $K=\sqrt F$.
Conversely, the union $K$ of an ascending sequence of radical ideals
$K_1\subset K_2\subset\cdots$ is radical: if a power of an element
lies in the union, it lies in one $K_i$, which then contains the element.
Write $K=\sqrt{(f_1,\ldots,f_n)}$. Each $f_j$ belongs to some
$K_i$, and a single $K_N$ contains all of them. Therefore
$K=\sqrt{(f_1,\ldots,f_n)}\subset K_N\subset K$, so the sequence
stabilizes. This proves ACC. Let
$$
\mathcal{F} = \{I \subset R \mid
\sqrt{I} = \sqrt{(f_1, \ldots, f_n)}\text{ for some }n
\text{ and }f_1, \ldots, f_n \in R\}.
$$
We now prove that $\mathcal F$ is an Oka family. First, $R\in\mathcal F$
using the generator $1$. For any $I$ and $a$, the product
$(I,a)(I:a)$ is contained in $I$. Conversely, for $x\in I$ we have
$x\in(I,a)$ and $x\in(I:a)$, since $xa\in I$; hence
$x^2\in(I,a)(I:a)$. These two containments prove the identity
$$
\sqrt{I} = \sqrt{(I, a)(I : a)}
$$
for every ideal $I\subset R$ and element $a\in R$.
If $\sqrt{(I,a)}=\sqrt{(f_1,\ldots,f_r)}$ and
$\sqrt{(I:a)}=\sqrt{(g_1,\ldots,g_s)}$, it follows that
$$
\sqrt I=\sqrt{(f_i g_j\mid 1\leq i\leq r,\ 1\leq j\leq s)}.
$$
Indeed, a prime contains a product of two ideals if and only if it
contains one of them, and replacing each ideal by another with the
same radical preserves precisely those primes. The radical is the
intersection of its containing primes by Lemma \ref{lemma-Zariski-topology}.
The displayed finite set of products therefore proves the Oka condition.
If the complement of $\mathcal F$ is nonempty, choose an ideal
in it as an upper bound for the empty chain. On the other hand,
if we have a nonempty totally ordered chain of ideals
$\{I_\alpha\}$ none of which are in $\mathcal{F}$, then the union
$I = \bigcup I_\alpha$ cannot be in $\mathcal{F}$ either. Otherwise
$\sqrt{I} = \sqrt{(f_1, \ldots, f_n)}$, then $f_i^e \in I$ for some $e$,
then $f_i^e \in I_\alpha$ for some $\alpha$ independent of $i$, then
$\sqrt{I_\alpha} = \sqrt{(f_1, \ldots, f_n)}$, contradiction.
Thus if the set of ideals not in $\mathcal{F}$ is nonempty, then it
has maximal elements and
exactly as in Lemma \ref{lemma-cohen} we conclude that $X$ is a
Noetherian topological space if and only if every prime ideal of $R$
is equal to $\sqrt{(f_1, \ldots, f_n)}$ for some $f_1, \ldots, f_n \in R$.
In detail, if every prime belongs to $\mathcal F$ but its complement
were nonempty, the chain argument and Zorn's lemma would give a maximal
excluded ideal, which is prime by Proposition \ref{proposition-oka},
a contradiction. Thus every ideal belongs to $\mathcal F$, giving
ACC for radical ideals by the equivalence just proved.
Conversely, ACC gives the finite radical presentation of every radical
ideal, in particular every prime. This proves the stated criterion
for the spectrum to be Noetherian.
\end{example}